\documentclass[../../../include/open-logic-section]{subfiles}

\begin{document}

\olfileid[id]{his}{set}{infinitesimals}
\olsection{Infinitesimal dan Diferensiasi}

Newton dan Leibniz menemukan kalkulus (secara independen) pada akhir abad
ke-17. Salah satu penerapan kalkulus yang sangat penting ialah
\emph{diferensiasi}. Secara kasar, diferensiasi bertujuan memberikan suatu
pengertian tentang ``laju perubahan'', atau gradien, suatu fungsi pada
suatu titik.

Berikut cara yang jelas untuk menggambarkan gagasan tersebut. Pertimbangkan
fungsi $f(x)=\nicefrac{x^2}{4}+\nicefrac{1}{2}$, yang digambarkan dengan
warna hitam di bawah ini:
\begin{center}
	\begin{tikzpicture}[scale=1]
	\draw[->, gray] (-1,0) -- (4.25,0) node[right] {$x$};
	\draw[->, gray] (0,-0.25) -- (0,5.25) node[above] {$f(x)$};

	\foreach \x/\xtext in {1/1, 2/2, 3/3, 4/4}
	\draw[shift={(\x,0)}, gray] (0pt,2pt) -- (0pt,-2pt) node[below] {$\xtext$};

	\foreach \y/\ytext in {1/1, 2/2, 3/3, 4/4, 5/5}
	\draw[shift={(0,\y)}, gray] (2pt,0pt) -- (-2pt,0pt) node[left] {$\ytext$};

	\draw[oldiagcolorC] (0.5, 9/16) -- (3.5, 57/16) -- (3.5, 9/16)--cycle;
	\draw[oldiagcolorD] (.5, 9/16) -- (2.5, 33/16) -- (2.5, 9/16) -- cycle;
	\draw[oldiagcolorE] (.5, 9/16) -- (1.5, 17/16) -- (1.5, 9/16) -- cycle;
	\draw[thick] (-1,.75) parabola bend (0,0.5) (4,4.5);
	\end{tikzpicture}
\end{center}
Andaikan kita hendak mencari gradien fungsi pada
$c=\nicefrac{1}{2}$. Kita mulai dengan menggambar suatu segitiga yang
hipotenusanya menghampiri gradien pada titik tersebut, misalnya segitiga
merah di atas. Jika $\beta$ merupakan panjang alas segitiga kita, tingginya
ialah $f(\nicefrac{1}{2}+\beta)-f(\nicefrac{1}{2})$, sehingga gradien
hipotenusanya ialah:
\[
\frac{f(\nicefrac{1}{2}+\beta) - f(\nicefrac{1}{2})}{\beta}.
\]
Jadi, gradien segitiga !!{colorC} kita, yang panjang alasnya~$3$, tepat
sama dengan~$1$. Hipotenusa segitiga yang lebih kecil, yaitu segitiga
!!{colorD} dengan panjang alas~$2$, memberikan hampiran yang lebih baik;
gradiennya ialah $\nicefrac{3}{4}$. Segitiga yang lebih kecil lagi, yakni
segitiga !!{colorE} dengan panjang alas~$1$, memberikan hampiran yang lebih
baik lagi; gradiennya ialah $\nicefrac{1}{2}$.

Segitiga yang semakin kecil memberikan hampiran yang semakin baik. Karena
itu, kita mungkin berkata seperti berikut: hipotenusa suatu segitiga dengan
panjang alas \emph{infinitesimal} memberikan gradien pada
$c=\nicefrac{1}{2}$ itu sendiri. Dengan cara ini, kita memperoleh suatu
rumus bagi turunan (pertama) fungsi $f$ pada titik $c$:
\[
{f'}(c) = \frac{f(c+\beta) - f(c)}{\beta}
\text{ dengan $\beta$ infinitesimal.}
\]
Secara kasar, inilah yang dikatakan Newton dan Leibniz.

Namun, karena mereka mengatakan demikian, kita harus bertanya: apakah
\emph{infinitesimal} itu? Timbul suatu dilema serius. Jika $\beta=0$, maka
$f'$ tidak terdefinisi dengan baik karena melibatkan pembagian oleh~$0$.
Namun, jika $\beta>0$, kita hanya memperoleh suatu \emph{hampiran} gradien,
bukan gradien itu sendiri.

Kekhawatiran ini bukan anakronisme. Berikut kritik Berkeley terhadap para
pengikut Newton:
\begin{quote}
Saya mengakui bahwa tanda dapat dibuat untuk menyatakan sesuatu atau tidak
menyatakan apa pun; dan karena itu, dalam notasi semula $c+\beta$, $\beta$
dapat menyatakan suatu pertambahan ataupun ketiadaan. Namun, apa pun yang
Anda jadikan maknanya, Anda harus bernalar secara konsisten dengan makna itu
dan tidak melanjutkan berdasarkan makna ganda; melakukan hal tersebut
merupakan sofisme yang nyata. (\citealt[\S{}XIII]{Berkeley1734},
variabel diubah agar sesuai dengan teks sebelumnya)
\end{quote}
Untuk membela kalkulus infinitesimal terhadap Berkeley, orang mungkin
menjawab bahwa pembicaraan tentang ``infinitesimal'' hanyalah kiasan.
Orang mungkin berkata bahwa, selama kita mengambil segitiga yang
\emph{sangat kecil}, kita akan memperoleh suatu hampiran garis singgung yang
\emph{cukup baik}. Berkeley juga mempunyai jawaban bagi hal ini: walaupun
hal itu mungkin cukup baik bagi rekayasa, ia merusak \emph{status}
matematika, sebab
\begin{quote}
kita diberi tahu bahwa \emph{in rebus mathematicis errores qu\`{a}m minimi
non sunt contemnendi}. [Dalam matematika, kesalahan yang paling kecil pun
tidak boleh diabaikan.] \citep[\S{}IX]{Berkeley1734}
\end{quote}
Petikan bercetak miring tersebut hampir merupakan kutipan kata demi kata
dari \emph{Quadrature of Curves} (1704) karya Newton sendiri.

Keberatan filosofis Berkeley sangat tajam. Meskipun demikian, kalkulus
merupakan usaha yang sangat berhasil, dan para matematikawan terus
menggunakannya tanpa terjerumus ke dalam kesalahan.

\end{document}
